
Results demonstrate that Pilot-Driven Viability Gates achieve 93.3\% accuracy for early identification of computationally infeasible hypotheses using synthetic validation. This section interprets findings, acknowledges limitations, and assesses broader implications.

\subsubsection{6.1 Key Findings}

\textbf{Perfect linear scaling is a synthetic artifact, but the mechanism is sound.}

The r=1.000 correlation between micro-pilot and full-scale overhead validates the framework's extrapolation logic in idealized conditions. However, this perfect linearity is unlikely in real experiments. Three sources of deviation are expected:

\begin{enumerate}
\item Measurement noise: Hardware variance, system load, I/O contention introduce timing fluctuations. Real correlation expected r=0.7-0.9.
\end{enumerate}

\begin{enumerate}
\item Non-linear scaling: Memory bottlenecks (100-sample dataset fits in cache, full dataset spills to RAM) or I/O overhead (data loading dominates for large datasets) create non-linear scaling patterns not captured by O\_full = k $\times O_{10}$.
\end{enumerate}

\begin{enumerate}
\item Type-specific scaling: Attention mechanisms (O($n^{2})$ in sequence length) may exhibit different k than gradient penalties (O(n)). CV=0.00\% across types is artifact; real data expected CV=10-30\%.
\end{enumerate}

Despite expected deviations, the r $\geq$ 0.7 threshold provides buffer. Even with measurement noise reducing correlation from r=1.000 to r=0.75, the mechanism remains valid. The risk is r < 0.7, where extrapolation breaks down---this failure mode should trigger framework recalibration (memory profiling for non-linear cases).

\textbf{Marginal Bayesian benefit reflects strong prior, not mechanism failure.}

Error reduction 40.91\% vs 40\% threshold (0.91 percentage point excess) might suggest fragile benefit. However, paired t-test (p=0.0003) confirms statistical significance, and absolute error reduction (84\%: 0.6966 $\rightarrow$ 0.1111) is substantial. The marginal relative reduction reflects synthetic data's strong prior (k=1.000 perfect scaling leaves little room for Bayesian refinement).

Real-world scenarios with weaker correlation (r=0.7-0.8) will have higher Gate 1 prior error, providing more opportunity for Gate 2 likelihood to refine predictions. Error reduction >50\% is hypothesized for real data. The mechanism works as designed; the marginal result is artifact of idealized conditions.

\textbf{High recall (96.2\%) on non-viable hypotheses minimizes false negatives.}

The framework's design priority---minimize missed non-viable hypotheses---is empirically validated. Only 1 of 26 non-viable hypotheses escaped Gate 1 (false negative rate 3.8\%). This false negative occurred at 10.3\% overhead, barely above the 10\% threshold, suggesting measurement uncertainty rather than systematic prediction failure.

The single false positive (viable hypothesis incorrectly stopped) occurred at 9.7\% overhead, also borderline. For non-borderline cases (overhead <9\% or >11\%), classification was 100\% accurate. This suggests incorporating confidence intervals: borderline cases (within $\pm 1$\% of threshold) should proceed to Gate 2 for refined prediction rather than stopping at Gate 1.

\subsubsection{6.2 Limitations}

\textbf{Synthetic Corpus Limits External Validity}

Validation used synthetic data with perfect linear scaling (k=1.000, r=1.000, CV=0.00\%). This establishes proof-of-concept---the framework's gate logic, statistical tests, and Bayesian updates work as designed. However, external validity is unknown: Does r $\geq$ 0.7 correlation hold on real published papers? Do real hypothesis types exhibit consistent scaling (CV < 30\%)?

Framework validated in synthetic contexts with perfect linear scaling. Real-world applicability contingent on real corpus validation showing r $\geq$ 0.7.

\textbf{Single Threshold Tested}

All validation used T=10\% overhead threshold (real-time deployment constraint). Framework behavior at permissive thresholds (T=50\% for offline batch processing, T=200\% for research contexts) is unknown. The 87\% non-viable prevalence at 10\% threshold may shift to 50\% at 50\% threshold or 10\% at 200\% threshold, changing accuracy characteristics.

Framework validated for strict deployment thresholds (10\%). Multi-threshold validation needed to establish threshold-accuracy curve.

\textbf{User Compliance Untested}

Gate 1 achieved 93.3\% classification accuracy (can identify non-viable hypotheses). The framework assumes researchers will act on stop signals (compliance behavior). This is untested. Potential failure modes include confirmation bias (ignoring Gate 1 stop signal for favored hypothesis), sunk cost fallacy (prior investment motivates full implementation despite stop signal), and publication pressure (negative results harder to publish).

Framework validated as predictive tool (93.3\% accuracy). Adoption as decision-making tool requires prospective user study measuring stop rate and time savings.

\textbf{Imbalanced Corpus May Inflate Accuracy}

With 26 non-viable and 4 viable hypotheses (87\% non-viable prevalence), the 93.3\% accuracy metric is dominated by high recall (96.2\%) on the large non-viable class. False positive rate (25\%: 1/4 viable stopped) is less representative due to small viable sample size.

Balanced corpus validation (15 viable, 15 non-viable at 50\% threshold) would test whether accuracy $\geq 80$\% holds with equal class distribution. Accuracy reduction to 85-90\% is anticipated under balanced prevalence but still above 80\% threshold.

Framework validated under imbalanced distribution (87\% non-viable). Balanced validation recommended to confirm accuracy $\geq 80$\% with 50-50 prevalence.

\subsubsection{6.3 Broader Impact}

\textbf{Positive Impacts:}

Reduced research waste: The h-e1 case shows 68.65\% overhead discovered after full implementation. Framework enables micro-pilot detection before wasted effort.

Systematic methodology: Fills gap in ML research workflow. Ablation studies test variations; complexity analysis provides theoretical bounds; viability gates add empirical early-stop protocol.

Quantified uncertainty: Bayesian posterior provides confidence intervals, not binary decisions. Researchers can assess borderline cases with probabilistic reasoning.

\textbf{Risks and Mitigation:}

False positives risk missed opportunities: 1 of 4 viable hypotheses incorrectly stopped (25\% FP rate on small sample). A viable hypothesis rejected at Gate 1 is a lost contribution. Mitigation: Borderline cases (confidence interval straddles threshold) proceed to Gate 2 for refined prediction. Only high-confidence non-viable predictions (O\_pred > 1.1T) trigger immediate stop.

Overreliance on micro-pilot: Researchers may skip hypothesis refinement, trusting Gate 1 predictions uncritically. Mitigation: Framework provides probability distributions and confidence intervals, encouraging critical evaluation.

Publication bias amplification: Stopped hypotheses yield no positive results, reducing negative result publication. Mitigation: Framework documentation encourages reporting Gate 1 stops as methodological contribution.

The framework assumes computational efficiency as valued constraint. In resource-constrained settings (limited compute budgets, energy efficiency mandates), this assumption aligns with sustainability goals. However, premature stopping may hinder exploratory research where ''inefficient'' hypotheses later inspire efficient variants. Balancing viability gates (feasibility-first) with exploratory investigation (curiosity-driven) is recommended rather than replacing the latter.

\subsubsection{6.4 Threats to Validity}

\textbf{Internal Validity:} Synthetic corpus provides controlled comparison (same hypotheses across all gates). No confounds from hardware variance, dataset differences, or implementation variations. Risk: Synthetic data may mask real-world complexities (memory bottlenecks, I/O overhead, non-linear scaling).

\textbf{External Validity:} The limiting threat. Synthetic validation establishes proof-of-concept but not generalizability. Real corpus validation is critical next step. Until r $\geq$ 0.7 confirmed on real published papers, claims scope to ''framework mechanics validated in synthetic conditions.''

\textbf{Construct Validity:} Accuracy, correlation, and error reduction measured correctly with standard statistical tests (binomial, Pearson, paired t-test). Time investment claim (''<1 hour'' for Gate 1) not empirically validated---synthetic runtime not representative of real micro-pilot cost. Prospective validation needed to confirm time savings claim.

\textbf{Statistical Conclusion Validity:} All results show strong statistical significance (p < 0.05). Effect sizes large for M1 (r=1.000) and M3 (93.3\% accuracy), marginal for M2 (40.91\% vs 40\%). Sample sizes adequate: 32 hypotheses (M1), 20 hypotheses (M2), 30 hypotheses (M3). Risk: Perfect correlation (r=1.000) and zero variance (CV=0.00\%) unlikely to replicate in real data.
